\section{Limitations}
\label{sec:limitations}

Every result in \cref{sec:results} is the output of a model. This section
states where the model is known to be inadequate, ordered by how much it would
change the conclusions.

\subsection{Recirculation}

\Cref{mod:disk} assumes quiescent inflow. By the exchange-time calculation of
\cref{eq:exchange} the machine passes a room's volume of air through its disks
every \SI{21}{\second}, so within a minute of a thirty-minute session the
assumption is false. The correction is not a coefficient; it requires solving
the room flow with a moving boundary. Two consequences are unquantified: the
degradation of rotor performance in its own recirculated wake, and the
persistent room-scale draught of order \SI{0.1}{\metre\per\second}. This is the
largest open item.

\subsection{Rotor aerodynamics and the flight power model}

The rotor model is momentum theory plus a single profile term plus three
near-field corrections. There is no blade element momentum theory, so the
radial distribution of inflow and loading is not resolved and the effects of
twist and taper enter only through lumped constants. There is no wake model, no
blade-vortex interaction, and no unsteady inflow: by \cref{as:quasisteady} the
inflow responds instantaneously, whereas its true time constant is of order
$R/\vind\approx\SI{0.1}{\second}$, comparable with the attitude loop. The
descent regime is excluded entirely, as noted after \cref{eq:axial}.

The flight simulation compounds this. Its power law \cref{eq:edot} is
calibrated at hover and has no inflow term, so it returns the same power for
a rotor at a given speed whether the vehicle is hovering or translating. The
axial correction derived in \cref{sec:momentum} is therefore not in the
mission energies of \cref{tab:flight,tab:coupled}. Those tables close the
energy balance around a near-hover approximation of the power; that is a
computational coupling, and it is exercised honestly, but it is not a
validation of manoeuvring energy. The figures are appropriate for indoor
following at walking speed and vertical rates of a few tenths of a metre per
second, and are not to be read as predictions for take-off, descent or
aggressive flight. Coupling an axial and oblique inflow model consistently to
both thrust and power is the next step, and it must be done to both at once:
a climb term in the power alone, against a hover-calibrated thrust, would be
inconsistent.

\subsection{Acoustics}

Three separate weaknesses. The source term \cref{eq:splmodel} is a broadband
scaling law fitted to four aircraft of one manufacturer; the exponents are
theoretically motivated but the anchor is empirical and the aircraft are all
conventional two-blade designs, so extrapolation to $\sigma=\rSolidity$
paddle blades is an extrapolation. The room model \cref{eq:roomlevel} uses a
single frequency-averaged absorption coefficient and assumes a diffuse field,
which fails below the Schroeder frequency \cref{eq:schroeder}, i.e.\ around
\SI{200}{\hertz}, which is exactly where the first blade harmonics lie. And
tonal content is computed but not summed into the reported level, so a design
with a strong \SI{\rBpf}{\hertz} fundamental and audible beating between rotors
could be unacceptable at a level the model calls fine. A measured prototype is
the only way to settle this.

\subsection{Uncertainty}

Every result is a single point estimate. The quantities that carry the most
uncertainty are the figure of merit of a low-Reynolds paddle rotor, the motor
mass constant, the screen drag coefficients, the electrical efficiency, the
cell specific energy at the discharge rate used, and above all the acoustic
anchor, whose calibration spread alone is \SI{1.4}{\decibel}. No propagation
of these uncertainties through the closure was performed, so the paper does
not show that the design ranking of \cref{tab:displays,tab:clearance} is
robust to them, nor that the candidate satisfies its constraints across their
plausible ranges. The margins reported are margins of a nominal model. A
proper treatment samples the parameter ranges, re-closes the design at each
sample, and reports the fraction of samples for which each constraint holds.

\subsection{Structures}

\Cref{mod:arm} covers bending stress, tip deflection and the first bending
mode of an idealised cantilever. It does not cover buckling, joints, bonded
interfaces, fatigue, or crash loads. For a machine intended to operate beside a
person, crash and containment analysis is not optional in practice, and none is
performed here. Guard rings are modelled only as a mass; their ability to
actually arrest a blade is asserted, not computed.

\subsection{Failure modes}

The controller of \cref{sec:control} has no failure mode: the allocation
\cref{eq:allocprob} assumes every rotor is available, and after the loss of one
the quadrotor described here falls. That is a statement about this controller
and not about quadrotors. Mueller and D'Andrea~\cite{mueller2014} show that a
quadrotor can be controlled after the loss of one, two or even three
propellers in a relaxed-hover mode in which the vehicle spins continuously
about a tilted axis. That capability is not implemented here, and a
continuously spinning \SI{1.7}{\kilo\gram} machine at arm's length from a
person is not thereby shown to be acceptable.

A hexacopter is not the answer either, or not simply. With the engine's
alternating-spin regular layout, removing one rotor and its opposite leaves an
allocation matrix of rank \rHexRank\ (\cref{sec:verification}): the four
remaining rotors cannot produce thrust, roll, pitch and yaw independently, so
the standard allocation loses a wrench direction, and lift arithmetic alone
would call for $\lambda\ge1.5$ to hover on four rotors of six, not the
$\lambda\ge2$ an earlier draft stated. Re-optimised under the same limits as
the quadrotor, the hexacopter (six rotors of \SI{\rHexD}{\milli\metre}) comes
out at \SI{\rHexGross}{\kilo\gram}, \SI{\rHexPower}{\watt} and
\SI{\rHexEar}{\decibel A} at the ear against \SI{\rQuadGross}{\kilo\gram},
\SI{\rQuadPower}{\watt} and \SI{\rQuadEar}{\decibel A} for the quadrotor
(\cref{tab:hex}), and its authority is lower on every axis:
$\rHexRoll/\rHexPitch/\rHexYaw$ against
$\rQuadRoll/\rQuadPitch/\rQuadYaw\ \si{\radian\per\second\squared}$. Its
weakest axis in angular acceleration is roll, not pitch: the roll torque
authority is the smaller of the two (\cref{tab:hex}, last column), because
with two rotors on the pitch axis the yaw constraint couples into the roll
differential, and the display slab makes $J_{xx}$ exceed $J_{yy}$. The row sums
$\sum_i|\mat M_{j+1,i}|$, which an earlier draft used and which predict pitch
as the weak axis, mislead here (\cref{app:proofs}). The failure case therefore
needs a different answer than more rotors, and none is offered here.

\begin{table}[t]
\centering
\caption{Quadrotor and hexacopter re-optimised under the same limits as
\cref{tab:displays}. Authorities by \cref{eq:authority}; the last column is
the ratio of roll to pitch torque authority.}
\label{tab:hex}
\footnotesize\setlength{\tabcolsep}{4pt}
\begin{tabular}{@{}lrrrrrrrrr@{}}
\toprule
Layout & $D$ & Gross & Power & At ear & Span & $\alpha_{\mathrm{roll}}$ & $\alpha_{\mathrm{pitch}}$ & $\alpha_{\mathrm{yaw}}$ & $\tau_{\mathrm{roll}}/\tau_{\mathrm{pitch}}$\\
\midrule
\inputrows{results/tab_hex_rows}
\end{tabular}
\end{table}

\subsection{Dynamic clearance}

\Cref{prop:governor} is a conditional guarantee, and each of its conditions
is a limitation. It is enforced at update instants on the governed reference,
not in continuous time and not on the vehicle. It transfers to the vehicle only
through the hypothesis $\sup_t\|\vect r-\vect r_g\|\le\varepsilon$, which
\cref{tab:flight} checks for the missions flown and which nothing proves for
missions not flown. It lapses whenever the caps and the keep-out cannot be
satisfied together, which the engine reports as a violation rather than
hides; on the reported runs this did not occur, but only because a mission
whose ideal position lies inside the sphere is started on the sphere rather
than inside it, a choice that removes the conflict from the start-up and not
from the product. And it protects a sphere of radius $\mathcal{R}$ about the
centre of mass; the rotor clearance reported alongside it is the physical
quantity, and the difference between the two is the conservatism of the
envelope, not a margin the design has earned.

\subsection{Estimation and autonomy}

\Cref{as:estimation} gives the controller exact knowledge of the state and of
the user's position. Real vision-based human tracking has latency, dropouts,
and occlusion, and the reference governor of \cref{mod:governor} is precisely
where those errors would surface: a keep-out sphere centred on an estimated
head position is only as safe as the estimate, and a latency of
\SI{100}{\milli\second} at a walking speed of \SI{1.4}{\metre\per\second} is
\SI{14}{\centi\metre} of sphere in the wrong place, more than the whole
tracking margin $\varepsilon$. No estimator is modelled and no sensor noise is
injected.

\subsection{Mission energy and the operating policy}

The windowed extrapolation $E_{\mathrm{mission}}=\bar P t_f$ of
\cref{sec:numerics} is exact only for a mission periodic within the window. For
a session with distinct phases it is an approximation of unquantified error,
and the battery of \cref{tab:coupled} is sized for the worst single phase
repeated for the whole session, which is a policy rather than a law. Reported
endurance figures also assume constant $\eb$; real cells lose capacity at high
discharge and low temperature, and the internal-resistance voltage sag that
would reduce available thrust at low state of charge is not modelled.

\subsection{Numerics}

The convergence check that matters for a sampled-data simulation, holding the
controller rate fixed and refining the integration substeps, was not run
(\cref{sec:numerics}). The differential evolution of \cref{eq:outerprog} is
stochastic and offers no optimality certificate: the reported designs are the
best found, not the best that exists, and the exterior penalty of
\cref{eq:penalty} means a returned point may violate a constraint slightly,
which \cref{tab:displays} reports as such. Runs are seeded for
reproducibility, which makes results repeatable without making them global.

\subsection{Scope}

Nothing was built. The paper computes what a set of models implies, and the
models were calibrated against published data for aircraft that are not this
aircraft. The most valuable next step is not a better model but a measured
rotor: a single \SI{\rDiam}{\milli\metre} two-blade paddle propeller on a
thrust stand in an anechoic space would settle \cref{tab:acoustic},
\cref{prop:fmblade} and \cref{eq:recorrection} simultaneously, and those three
carry most of the uncertainty in the result.
